\chapter{INTRODUÇÃO}

A Universidade Tecnológica Federal do Paraná (UTFPR), instituição pública federal orientada pelos princípios da educação tecnológica, reafirma em seu Plano de Desenvolvimento Institucional \cite{UTFPR_PDI_2023} o compromisso com a formação cidadã, o desenvolvimento científico e a inovação tecnológica, em permanente diálogo com as demandas sociais, produtivas e ambientais. Fundamentada na integração entre ensino, pesquisa, extensão e inovação, a UTFPR consolida sua identidade como universidade multicampi comprometida com o desenvolvimento sustentável, a ampliação do acesso à educação superior de qualidade e a construção de soluções para os desafios contemporâneos da ciência, da tecnologia e da transformação digital.

Nesse contexto institucional insere-se o Projeto Pedagógico do Curso (PPC) de Bacharelado em Ciência de Dados e Inteligência Artificial (CDIA), vinculado ao Departamento Acadêmico de Engenharia Elétrica e Computação (DAEEC) do campus Londrina. O curso fundamenta-se nos valores institucionais da UTFPR — ética, desenvolvimento humano, integração social, inovação, qualidade, excelência e sustentabilidade — e contribui diretamente para o cumprimento de sua missão de desenvolver a educação tecnológica de excelência por meio do ensino, da pesquisa e da extensão, promovendo uma interação ética, produtiva, sustentável e inovadora com a sociedade \cite{UTFPR_PDI_2023}.

A estrutura multicampi da UTFPR, presente em treze campi distribuídos pelo estado do Paraná, constitui um arranjo institucional singular que articula vocações regionais, parcerias com o setor produtivo, parques tecnológicos, incubadoras e ecossistemas locais de inovação. Essa capilaridade fortalece a inserção regional da universidade e potencializa sua atuação frente às demandas nacionais e internacionais em ciência, tecnologia e inovação. No âmbito desse planejamento estratégico, a criação do curso de CDIA no campus Londrina está prevista no PDI 2023–2027 \cite{UTFPR_PDI_2023} e alinha-se ao Programa Universidades Inovadoras, promovido pelo Ministério da Educação, que orienta a modernização da educação superior frente às transformações científicas, tecnológicas e produtivas em curso.

O curso integra o conjunto de novos cursos concebidos de forma articulada entre os campi da UTFPR, segundo o modelo institucional UTFPR (STEM$^{IA}$), orientado às áreas de Ciência, Tecnologia, Engenharia e Matemática (STEM), potencializadas pelo uso ético, crítico e estratégico da Inteligência Artificial. O avanço acelerado de tecnologias como ciência de dados, inteligência artificial, automação inteligente, sistemas ciberfísicos e soluções digitais têm transformado profundamente os processos produtivos, os ambientes organizacionais e os modos de trabalho, demandando novos perfis profissionais com sólida formação científica, competências técnicas integradas e responsabilidade ética e social.

O desenvolvimento deste projeto pedagógico considera as legislações e normativas nacionais e institucionais vigentes \cite{LDB_1996, UTFPR_RODP_2019}, bem como diretrizes técnicas e científicas da área, além das especificidades e demandas locais e regionais, conforme estabelecido na Lei nº 9.394/96 (Lei de Diretrizes e Bases da Educação Nacional) \cite{LDB_1996}. O PPC explicita as ações educativas, os princípios formativos e as características estruturantes do curso, constituindo-se como instrumento orientador das práticas acadêmicas e pedagógicas dos profissionais a ele vinculados.

Nesse cenário, destaca-se o município de Londrina como um importante polo regional de Tecnologia da Informação e Comunicação (TIC). A cidade abriga um expressivo número de empresas de base tecnológica, gera milhares de empregos no setor e apresenta forte integração entre tecnologia, agronegócio, indústria, saúde, comércio e serviços. Esse ecossistema inovador cria condições favoráveis para a formação de profissionais altamente qualificados, capazes de aplicar técnicas de Ciência de Dados e Inteligência Artificial na solução de problemas complexos e no desenvolvimento de soluções tecnológicas de alto impacto econômico e social.

A implantação do curso de CDIA no campus Londrina responde diretamente a esse contexto regional, ampliando a oferta de formação pública, gratuita e de excelência em uma área estratégica para o desenvolvimento local e regional. O curso contribuirá para o fortalecimento do ecossistema de inovação da região, ao formar profissionais aptos a atuar em empresas de tecnologia, organizações públicas e privadas, centros de pesquisa e iniciativas empreendedoras, promovendo a transferência de conhecimento, a inovação aplicada e o desenvolvimento sustentável.

Adicionalmente, as unidades curriculares do curso foram estruturadas de modo a dialogar com os Objetivos de Desenvolvimento Sustentável (ODS) da Agenda 2030 da Organização das Nações Unidas, reforçando o compromisso institucional da UTFPR com a formação de profissionais capazes de enfrentar desafios globais de natureza tecnológica, social, ambiental e ética. Destacam-se, nesse sentido, os ODS relacionados à educação de qualidade, ao trabalho decente e crescimento econômico, à indústria, inovação e infraestrutura e à redução das desigualdades.

Ao integrar rigor científico, prática profissional, inovação e responsabilidade social, o curso de CDIA do campus Londrina reafirma a missão social da universidade pública e fortalece o papel da UTFPR como protagonista na formação de profissionais capazes de liderar processos de desenvolvimento tecnológico e inovação com impacto regional e nacional. Dessa forma, a proposta de implantação do curso consolida-se como uma iniciativa estratégica para o fortalecimento do setor tecnológico, da economia regional e do desenvolvimento sustentável do Paraná e do Brasil.

\section{Oferta do Curso}

O curso de CDIA será ofertado no período noturno para melhor atender ao público-alvo ao permitir que os discentes possam realizar estágios ou trabalhar durante o dia. A modalidade do curso é presencial com parte da carga horária em Ensino a Distância (EaD), respeitando o limite máximo de 30\% da carga horária total para cursos presenciais.

O tempo previsto para integralização da carga horária total do curso de 3.210 horas é de 8 semestres.

Serão ofertadas 80 vagas anuais, com entradas semestrais de 40 vagas em cada semestre. A oferta de disciplinas aos discentes também será realizada no regime semestral para melhor se adequar ao ingresso de estudantes.